\section{The Witnesses after 1846}

The children did not disappear into a devotional tableau. They grew into adults beneath attention they had not sought and institutions that did not always know how to protect them. Their later lives contain courage, instability, faith, conflict, ordinary failure, manipulation by others, and competing posthumous portraits. A serious account neither canonizes them by association nor destroys them to simplify a controversy.

\subsection{The Curé d'Ars episode}

On 25 September 1850 Maximin met Jean-Marie Vianney at Ars. The journey was promoted by adults despite episcopal concern and was entangled with Maximin's vocational questions and the hopes of political enthusiasts who wanted the secret to support their causes. The conversation left Vianney believing that Maximin had admitted invention or lying. The saint communicated his doubt to others.

Maximin denied retracting his apparition testimony. In statements on 2 November he said Vianney had not asked him directly about the apparition or secret as later alleged, and he explained that difficulty hearing or understanding had led him to answer yes or no without grasping questions. Other accounts place a skeptical vicar and several exchanges around the private interview. There is no verbatim transcript.

The episode is genuinely adverse and cannot be erased by saying ``a saint later believed.'' A discerning reader asks why Vianney understood a retraction and why Maximin's explanations took the form they did. Yet the same record does not permit conversion of a disputed, privately remembered exchange into a certain confession. De Bruillard knew of the controversy and still judged the event positively in 1851 after further inquiry.

Later reports say Vianney regained confidence in La Salette after prayer and a providential incident; their details and chronology vary. Neither a saint's doubt nor his reported reassurance is the competent act. Jean Stern's study \emph{Le curé d'Ars et le message authentique de La Salette}, with a preface by Bishop Guy de Kerimel, treats the confusion as a lesson in returning to the public message and refusing secret-centered distortions.

\sourceboundary{Maximin's statements, Vianney's reported recollections, associated correspondence, and Stern's documentary analysis.}{They establish a serious 1850 misunderstanding or contradiction that affected reception and was known before the judgment.}{They do not furnish an agreed transcript, prove either man's bad faith, settle the apparition by sanctity, or authenticate a later secret.}

\subsection{Maximin: a witness without a stable public role}

After the event Maximin was educated in several institutions and explored a priestly vocation without proceeding to ordination. His later path included studies, brief employment, an unsuccessful attempt at medicine, pharmacy-related work, travel, and dependence on patrons. In 1865 he enlisted for six months in the Papal Zouaves, first serving in a medical capacity and then as a soldier; his unit saw no major campaign that year. Political and devotional circles repeatedly tried to make his secret serve monarchist or apocalyptic programs.

Maximin could be impulsive and sometimes answered intrusive questioners flippantly. He was also subjected from childhood to thousands of questions and to adults who offered money, careers, travel, or ideological importance. Psychological reconstruction at a distance cannot replace evidence. His uneven life makes simplistic claims of a permanently serene ``holy seer'' untenable, but poverty and vocational failure are not evidence of fraud.

In 1866 he published \emph{Ma profession de foi sur l'apparition de Notre-Dame de La Salette}, answering attacks and reaffirming the event. Later editions reproduced de Bruillard's 1851 act. Texts and predictions also circulated under or around his name, and he rejected at least some later attributions. The safe corpus remains his early testimony, the sealed 1851 manuscript, and writings with demonstrable authorship.

Maximin died at Corps on 1 March 1875, aged thirty-nine. The Missionaries' official chronology records the date. He died within the Church and continued to affirm the apparition. That perseverance matters as witness evidence; it is not a canonization, proof of every story, or fulfillment of his secret.

\subsection{Mélanie: vocation, movement, and expansion}

Mélanie entered the care of the Sisters of Providence at Corenc and used the name Sister Marie of the Cross. She desired religious profession. Ginoulhiac did not permit her to make final profession there, judging her insufficiently mature and seeking distance from further visionary elaboration. She then moved through religious settings in England, France, Greece, and Italy, at times under temporary vows or structured obedience but without one stable lifelong institute.

Her mobility has been narrated either as persecution of a saint or proof of pathological instability. The documentary reality is less convenient. Some clergy and religious testified to piety, poverty, courage, and charity; others found her visionary claims, independence, political associations, or plans for a new order imprudent. She experienced genuine vulnerability and depended on patrons. She also exercised agency and persisted in interpretations that competent authorities restricted.

From the 1850s onward she produced or authorized increasingly expansive secret texts and claimed a mission concerning apostles of the last times and a religious rule. She lived for years at Castellammare and received support from southern Italian churchmen, including Bishop Zola. The 1879 booklet gave her expanded account its durable printed form. She continued to defend it despite Roman withdrawal instructions, while admirers produced commentaries and political applications.

Mélanie died at Altamura in the night of 14--15 December 1904. Local devotion surrounds her memory, and advocates have collected testimony to extraordinary holiness. She has not been beatified or canonized. Titles such as ``Blessed Mélanie'' or the use of alleged childhood visions and stigmata as established facts are not warranted by a recognized canonization judgment.

\subsection{Later character cannot bear the whole case}

Ginoulhiac's 1854 principle remains wise. An event in 1846 is not retroactively caused or uncaused by everything a witness does in 1879. A credible witness can later err, interpret, exaggerate, or become attached to a mistaken mission. A morally imperfect person can tell the truth about an earlier fact. Conversely, personal piety does not guarantee every recollection or prophecy.

The later lives remain relevant in proportion:

\begin{itemize}
  \item sustained affirmation bears on sincerity, but not directly on supernatural cause;
  \item changing texts bear on confidence in exact wording and later claims;
  \item responses to legitimate authority bear on prudence, but require exact orders and circumstances;
  \item financial or political exploitation can explain distortions, but must be proved rather than insinuated;
  \item spiritual fruit in others may support pastoral reception, but does not validate every source layer.
\end{itemize}

The Church's mission begins precisely because the witnesses cannot remain the permanent managers of a worldwide devotion. The public message enters preaching, sacrament, scholarship, and ecclesial oversight. The children are honored by telling the truth about their humanity, not by forcing them to carry an impossible charism of lifelong inerrancy.

\subsection{Two opposite abuses of the witnesses}

\begin{threestable}{Abuse}{What it does}{Corrective}
Devotional idealization & Treats poverty, lack of education, or later suffering as proof of sanctity and every text as Marian dictation. & Respect the children, date every source, and reserve judgments of heroic virtue to competent authority.\\
Polemical destruction & Collects rumors, failed careers, or conflict as if one flaw disproved every earlier statement. & Test each claim, distinguish sincerity from causation, and remember the bishop judged with major controversies already known.\\
Political capture & Makes a child's secret endorse a dynasty, party, nation, anti-clerical campaign, or theory of papal legitimacy. & No private revelation supplies Catholic political allegiance; expired and elastic predictions receive no privileged authority.\\
\end{threestable}

Mélanie's expanded accusations against clergy require an additional safeguard. The Church must confront clerical sin and abuse with truth, justice, safeguarding, and care for victims. A sweeping private-revelation denunciation is neither evidence against a particular person nor permission to distrust the sacramental priesthood. Conversely, invoking the public message's approval can never silence credible allegations. Reconciliation depends on truth.
